\SourceLeaf{272}
\Article{5}
\indent Ich begleite die Vorschriften dieser ganzen Abhandlung mit einer auf das schärfste durchgeführten numerischen Anwendung, welche andern, die zur Verarbeitung ihrer Messungen die hier vorgetragene Methode benutzen wollen, entweder als Rechnungsmuster zur Construction der erforderlichen Hülfstafeln, oder auch schon unmittelbar als Hülfsapparat für einen grossen Theil der gemässigten Zone dienen kann. In den meisten Fällen wird man übrigens sich mit einer \textit{viel} geringern Schärfe begnügen können.

Als Normalbreite wähle ich \(52^0\ 40'\), welche ungefähr dem mittlern Parallelkreise des Königreichs Hannover entspricht; da es jedoch in einigen Beziehungen vortheilhafter ist, wenn für die Normalbreite auf der Kugel, als wenn für die auf dem Ellipsoid eine runde Zahl gewählt wird, so setze ich
\[
Q=52^0\ 40'\ 0''
\]
\indent Die Rechnung führe ich nach den neuesten von \textsc{Bessel} aus den Gradmessungen abgeleiteten Erddimensionen (Astronomische Nachrichten 19. Band S. 116), wonach, die Toise zur Einheit angenommen,
\[
\begin{aligned}
\log a&=6{,}5148235337\\
\log\cos\varphi&=9{,}9985458202
\end{aligned}
\]
Es folgt hieraus, mit Hülfe der zehnziffrigen Logarithmen,
\[
\begin{aligned}
\varphi&=4^0\ 41'\ 9''98262\\
\log e&=8{,}9122052079\\
\zeta&=1^0\ 43'\ 26''80402\\
\eta&=2\quad15\quad42\ 34083\\
P&=52\quad42\quad2{,}53251\\
\log\alpha&=0{,}0001966553\\
\theta&=3^0\ 43'\ 34''24669\\
\log\frac{1}{k}&=0{,}0016708804\\
\log A&=6{,}5152074703
\end{aligned}
\]
Nimmt man das französische gesetzliche Meter als Einheit an, so wird
\[
\log A=6{,}8050274003
\]
